\documentclass[a4paper,12pt]{article}
\usepackage[utf8]{inputenc}
\usepackage{amsmath}
\usepackage{graphicx}
\usepackage{hyperref}
\usepackage{geometry}
\geometry{left=3cm, right=2.5cm, top=2.5cm, bottom=2.5cm}

\title{\textbf{Projeto Final Gitflow}}
\author{Aníbal Freire}
\date{\today}

\begin{document}

% Capa com imagem
\begin{titlepage}
    \centering
    \includegraphics[width=0.5\textwidth]{istec.jpeg} % Insira o caminho da imagem aqui
    \vspace{1cm}

    {\LARGE \textbf{ISTEC - Instituto Superior de Tecnologias Avançadas do Porto} \par}
    \vspace{1.5cm}

    {\Large \textbf{Projeto Final Gitflow} \par}
    \vspace{0.5cm}

    {\large DAS - Desenvolvimento Ágil de Software \par}
    \vspace{5cm}

    \textbf{Aníbal Freire} \\
    Curso: Desenvolvimento de Software \\

    \vfill

    {\large \today\par}
\end{titlepage}

\newpage
\tableofcontents
\newpage

% Conteúdo do relatório
\section{Título do Projeto}
Projeto Final Gitflow.

\section{Introdução}

O objetivo deste projeto é criar um repositório no GitHub, usando o GitFlow, no qual devemos implementar boas práticas de controle de versão e gestão de branches.

\section{Membros da Equipa e Tarefas}

Aníbal Freire
\begin{itemize}
    \item Criação do Github
    \item Criação do relatórii
\end{itemize}
\vspace{0.2cm}

\noindent
Diogo Lima
\begin{itemize}
    \item Auxílio com pesquisa
    \item Realizaçºão de testes das permissões do Githab
\end{itemize}
\vspace{0.2cm}

\newpage
\section{Link para o repositório}

\url{https://github.com/zewo2/das-final}

\section{Código git usado}

Para a criação do repositório local e uso do git flow para a manipulação e gestão deste foram usados os seguintes comandos em git bash:

\vspace{0.3cm}

\noindent\fbox{\begin{minipage}{\dimexpr\textwidth-2\fboxsep-2\fboxrule\relax}
\centering

git clone https://github.com/zewo2/das-final

cd das-final

git flow init

\end{minipage}}

\vspace{0.3cm}

Depois, para a subsequente manipulação do repositório remoto, upload das alterações e controlo de versão foram usados os seguintes comandos:

\vspace{0.3cm}

\noindent\fbox{\begin{minipage}{\dimexpr\textwidth-2\fboxsep-2\fboxrule\relax}
\centering

git remote add origin https://github.com/zewo2/das-final

\end{minipage}}

\vspace{0.3cm}

Após a adição da origem remota, foi realizada a base inicial do relatório e colocada no repositório local, sendo estas alterações depois enviadas para o remoto:

\vspace{0.3cm}

\noindent\fbox{\begin{minipage}{\dimexpr\textwidth-2\fboxsep-2\fboxrule\relax}
\centering

git status

git add .

git commit -m "mensagem de commit"

git push origin develop

\end{minipage}}

\vspace{0.3cm}

Após concluir o relatório e fazer as cinco (5) alterações e commits necessários, foi criada uma release que foi publicada para motivos de arquivo:

\vspace{0.3cm}

\noindent\fbox{\begin{minipage}{\dimexpr\textwidth-2\fboxsep-2\fboxrule\relax}
\centering

git flow release start Rel-1.0 develop

git flow release publish Rel-1.0

git flow release finish Rel-1.0 -m "v1" -k

\end{minipage}}

\vspace{0.3cm}

Como foram denotados alguns erros no relatório, foram usados os comandos referentes às funções de hotfix para corrigir os mesmos:

\vspace{0.3cm}

\noindent\fbox{\begin{minipage}{\dimexpr\textwidth-2\fboxsep-2\fboxrule\relax}
\centering

git flow hotfix start Rel-1.1 develop

git flow hotifx publish Rel-1.1

git flow hotfix finish Rel-1.1 -k

\end{minipage}}

\section{Ferramentas e Tecnologias}
Para realizar este projeto usufrui das seguintes ferramentas:
\begin{itemize}
    \item Overleaf - Para escrever este relatório
    \item Git - Controlar repositórios locais e remotos
    \item MinGW - Para interagir com o Git.
    \item Github - Para  controlar a versão final e certificar que esta estava sincronizada através de todos os dispositivos e permitir a colaboração remota com o resto da equipa.
    \item GIMP - Para editar algumas imagens usadas no relatório
\end{itemize}

\subsection{Overleaf}
    \includegraphics[width=0.2\textwidth]{overleaf.jpg} % Trocar imagem com logo do Overleaf
    \vspace{0.5cm}

O Overleaf [1] é um editor de Latex que não requer qualquer instalação, por ser online, cujo permite a fácil colaboração entre vários membros de um projeto.

Este foi projetado para tornar o processo de escrita e criação de relatórios técnicos mais fácil e rápido, através da utilização de templates e dos seus intuitivos editores Visual e de Código, permitindo utilizadores com diferentes competências realizarem o seu projeto de forma intuitiva.

O Overleaf permite ainda a utilização de gestores bibliográficos para a inserção de citações e referências bibliográficas, sendo alguns exemplos deste tipo de ferramentas o Zotero e o Mendeley.

No âmbito deste projeto, utilizei o Overleaf para me permitir escrever este relatório.

\subsection{Git}
    \includegraphics[width=0.35\textwidth]{git.png} % Trocar imagem com logo do Git
    \vspace{0.5cm}

O Git [2] é um sistema de controle de versões distribuído, usado principalmente no desenvolvimento de software, mas pode ser usado para registar o histórico de edições de qualquer tipo de arquivo (Exemplo: alguns livros digitais são disponibilizados no GitHub e escrito aos poucos publicamente).

O Git foi inicialmente projetado e desenvolvido por Linus Torvalds para o desenvolvimento do kernel Linux, mas foi adotado por muitos outros projetos.

Cada diretório de trabalho do Git é um repositório com um histórico completo e habilidade total de acompanhamento das revisões, não dependente de acesso a uma rede ou a um servidor central.

No âmbito deste projeto, utilizamos o Git para trabalhar com repositório locais e remotos com sincronização e controlo de versão.

\subsection{MinGW}
    \includegraphics[width=0.35\textwidth]{mingw.png} % Trocar imagem com logo do MinGW
    \vspace{0.5cm}

O MinGW [3] (Minimalist GNU for Windows) é uma versão portada para Microsoft Windows do conjunto de ferramentas GNU. 

Este software inclui um conjunto de arquivos cabeçalho para a API do Windows que permite aos programadores usarem o GCC para criarem programas nativos em Windows sem precisarem contar com uma emulação em tempo real de um sistema Unix-like.

Este era originalmente uma ramificação do Cygwin, que fornece um suporte Unix-like maior para Windows. O MinGW é software livre, estando os seus cabeçalhos Win32 disponíveis em domínio público enquanto as ferramentas estão em GNU General Public License.

No âmbito deste projeto, utilizamos o MinGW como interface para poder interagir com o Git.

\subsection{Github}
    \includegraphics[width=0.25\textwidth]{github.png} % Trocar imagem com logo do github
    \vspace{0.5cm}

O GitHub [4] é uma plataforma de hospedagem de código-fonte e arquivos com controlo de versão utilizado o Git. Ele permite que programadores, utilitários ou qualquer utilizador registado na plataforma contribuam em projetos privados e/ou Open Source de qualquer lugar do mundo. 

Este é amplamente utilizado por programadores para a divulgação dos seus trabalhos ou para que outros programadores contribuam para o seu projeto, para além de promover fácil comunicação através de recursos que relatam problemas ou misturam repositórios remotos (issues, pull request).

O GitHub é mundialmente usado e chega a ter mais de 36 milhões de utilizadores ativos mundialmente contribuindo em projetos comerciais ou pessoais. Hoje o GitHub abriga mais de 100 milhões de projetos, alguns deles sendo conhecidos mundialmente, como o WordPress, GNU/Linux, Atom e Electron.

No âmbito deste projeto, usamos o Github para controlar as versões atuais do projeto, podermos colaborar remotamente e corrigir erros sem haver necessidade de processos complexos e demorados.

\subsection{GIMP}
    \includegraphics[width=0.3\textwidth]{gimp.png} % Trocar imagem com logo do GIMP
    \vspace{0.5cm}

O GIMP (GNU Image Manipulation Program) (originalmente (General Image Manipulation Program) [5], ou traduzido em Português (Programa de Manipulação de Imagem do GNU) é um programa de software "open-source" focado na criação e edição de imagens bitmap, e desenhos vetoriais.

O GIMP foi criado por estudantes como uma alternativa livre e gratuita ao Adobe Photoshop. Foi um projeto universitário que amadureceu bastante e hoje alcança expressiva popularidade, sendo utilizado por artistas amadores e profissionais do ramo.

O GIMP foi criado para a criação ou manipulação de imagens e fotografias, podendo também criar gráficos, logótipos, redimensionar fotos, alterar cores, combinar imagens utilizando camadas, remover partes indesejadas e converter arquivos entre diferentes formatos de imagem digital.

No âmbito deste projeto, utilizei o GIMP para redimensionar e alterar algumas imagens.

\section{Webografia}
\begin{itemize}
    \item [1] Overleaf, "https://overleaf.com/about/features-overview", 2025

    \item [2] Git, "https://pt.wikipedia.org/wiki/Git", 2025

    \item [3] MinGW, "https://pt.wikipedia.org/wiki/MinGW", 2025
    
    \item [4] GITHUB, "https://docs.github.com/pt/get-started/
    \newline
    using-git/about-git", 2025
    
    \item [5] GIMP, "https://docs.gimp.org/2.8/pt-BR/introduction.html", 2025
    
\end{itemize}

\section{Considerações Finais}

Conclusão

\end{document}